\chapter{领域知识组织}

\section{检索与断言}

检索采用纯标准库实现的 BM25，中文按字符二元组切分。
刻意不使用分词与向量检索库，理由是要保证评委在自己的机器上克隆下来即可运行，
少一个依赖就少一个现场翻车的可能。后续升级只需替换检索接口的实现，调用方不变。

生成智能体输出的最小单元是断言，每条必须可独立判断真伪，且携带知识库切片编号。

\section{溯源可信度：压倒其余全部指标的一件事}

\begin{keynote}{知识库本身没核实，评测数就不是效果数}
整套幻觉检测的逻辑是“断言必须被知识库支撑”。
那么反过来，\emph{知识库本身如果是错的，审核闸只会把错误认证为正确}。
跑分中的“幻觉率百分之零”，真实含义是“生成内容与知识库一致”，
\emph{不是“内容正确”}。这两句话差得很远，而对外材料极容易把前者说成后者。
\end{keynote}

骨架版交付时，二十六条切片中有二十一条的出处是编造的占位文本，
形如“《工业机器人操作与运维》第 4 章”，看上去毫无破绽。
这比明显的胡编更危险，因为它会让读者、评委、甚至团队自己误以为内容有据可查。

现已联网核实，结果见表 \ref{tab:verify}。

\begin{table}[htbp]\centering\small
\caption{知识库核实结果}\label{tab:verify}
\begin{tabular}{@{}p{2.8cm}p{2.2cm}p{6.6cm}@{}}
\toprule
条目 & 结论 & 说明 \\
\midrule
四个报警代码 & 内容正确 & SRVO-001、002、005、062 分别对应操作面板急停、
示教器急停、机器人超程、脉冲编码器备用电池电压不足，与多份厂商资料一致 \\
示教限速 & 内容正确 & 示教时运动速度应低于 250\unit{mm/s}\cite{ansafety} \\
围栏高度与安全距离 & \emph{出处错配} & 原稿把数值挂在一条未经核实的国标条款名下。
内容或许合理，但挂到具体标准上属于错配 \\
其余二十一条 & 未核实 & 含标定公差、润滑周期、子程序层数上限等 \\
\bottomrule
\end{tabular}
\end{table}

核实过程中出现了一个值得记住的现象：某维修服务商网站把 SRVO-001 说成
“伺服放大器过载”，与多份官方手册摘录完全矛盾。
\emph{单一网络来源不足以定案}，这不是理论上的顾虑。

\section{核实状态的五档取值}

每条切片现在都带核实状态字段，取值分五档，见表 \ref{tab:fcstatus}。

\begin{table}[htbp]\centering\small
\caption{核实状态取值与写入权限}\label{tab:fcstatus}
\begin{tabular}{@{}p{3cm}p{5.6cm}p{2.6cm}@{}}
\toprule
状态 & 含义 & 谁能写 \\
\midrule
\cd{unverified} & 未核实（默认） & 机器 \\
\cd{sourced} & 来自真实文档，可回溯到文件与位置 & 机器 \\
\cd{machine\_checked} & 找到外部证据且多来源一致 & 机器 \\
\cd{disputed} & 证据不足、来源单一或多来源矛盾 & 机器 \\
\cd{refuted} & 找到明确矛盾的外部证据 & 机器 \\
\cd{verified} & \emph{人工确认无误} & \emph{只有人} \\
\bottomrule
\end{tabular}
\end{table}

最后一行是整个设计的支点。一旦机器能写 \cd{verified}，
整条溯源链就退化为“模型说它查过了”，与不核实无异；
写回函数中有断言直接拦住此类写入。

\cd{sourced} 与 \cd{verified} 的区别也需要讲清楚。
从真实文档切出来的内容是“有出处的”，不等于“经过核实的”——
出处真实只说明这句话确实印在那本书上，不说明那本书是对的。
但 \cd{sourced} 本身已是巨大进步：它的出处是文件名加页码加文件指纹，
任何人都能翻回原文核对，而骨架版那二十一条是查无此书。

\section{能不能让模型自动核实}

能，但有一条铁律，破了就前功尽弃：
\emph{模型的判断本身永远不能成为核实依据，核实依据只能是外部证据。}

直觉做法是把切片丢给大模型问“这条对不对”，模型说对就标成已核实。
\emph{这么做比不核实还糟。} 知识库本来就是大模型写出来的，
再让大模型判它对不对，等于让同一个来源自己给自己背书——
模型倾向于认可符合自己先验的说法，而那些说法恰恰是它自己生成的。
结果是把“未核实”洗成“已核实”，标签变了、可信度一点没变，
错误从此戴上一顶合规的帽子，比裸奔更难被发现。

核实智能体按这条铁律实现。模型在其中只干两件事：
把切片改写成检索词；读取\emph{检索回来的外部资料}判断是否支持该切片。
第二件事仍是判断，但判断对象是别人写的资料，而非自己的记忆。
检索不到资料就判存疑，\emph{绝不允许回落到模型记忆}。

判定遵循三条规则。\emph{多来源}：至少两个互相独立的域名一致才算通过，
同一域名下的多个页面不构成独立来源。
\emph{矛盾优先}：只要有一个来源给出矛盾结论就判被推翻，哪怕支持方更多；
测试中构造过三比一的场景，矛盾方胜出。
\emph{无证据不放行}：检索源没接上时全部判存疑，这是正确行为而非故障。

矛盾优先这一条对应本项目真实踩过的坑：上一节提到的 SRVO-001 争议，
当时是人发现的，现在系统能自动把它顶出来并附上证据清单。

\begin{keynote}{这个工具是省人力的，不是替人拍板的}
它把证据摆到人面前并排好优先级，最后那一下还得人来点——
检索到的网页本身也可能是错的。
\end{keynote}

\bibliographystyle{gbt7714-2005}
\bibliography{refs}
